\addcontentsline{toc}{chapter}{KATA PENGANTAR}
\begin{center}
    \textbf{KATA PENGANTAR}
\end{center}

\vspace{1cm}

\begin{center}
    \textit{Assalamu'alaikum warahmatullahi wabarakatuh.}
\end{center}

\vspace{0.5cm}

\hspace*{1cm}Puji syukur penulis panjatkan ke hadirat Allah SWT atas segala rahmat, hidayah, serta kekuatan yang telah diberikan sehingga penulis dapat menyelesaikan skripsi ini. Shalawat serta salam semoga selalu tercurah kepada junjungan kita, Nabi Muhammad SAW, beserta keluarga, sahabat, dan seluruh pengikutnya hingga akhir zaman.

\hspace*{1cm}Skripsi ini disusun sebagai salah satu syarat untuk memperoleh gelar Sarjana Pendidikan pada Program Studi Pendidikan Matematika, Universitas Ahmad Dahlan. Penulis menyadari bahwa proses penyelesaian skripsi ini tidaklah mudah. Berbagai rintangan dan tantangan telah dihadapi, namun berkat kerja keras, kesabaran, serta dukungan dari berbagai pihak, skripsi ini akhirnya dapat diselesaikan. Penulis juga menyadari bahwa skripsi ini baru dapat selesai pada semester ke-14, yang merupakan perjalanan panjang dalam menempuh pendidikan di perguruan tinggi.

\hspace*{1cm}Skripsi ini penulis persembahkan untuk ibu tercinta yang telah meninggal dunia pada tanggal 28 Agustus 2025. Semoga ilmu yang penulis peroleh selama berkuliah dapat menjadi gambaran bagaimana beliau mendidik penulis dengan penuh kasih sayang selama 25 tahun. Ibu adalah sumber motivasi terbesar bagi penulis untuk terus berjuang menyelesaikan skripsi ini. Kepada ayah yang selalu menjadi teladan dan bekerja keras membiayai pendidikan anak-anaknya, penulis mengucapkan terima kasih yang tak terhingga. Beliau adalah sosok yang tidak pernah lelah mendukung dan mendoakan penulis.

\hspace*{1cm}Penulis juga menyampaikan rasa terima kasih yang sebesar-besarnya kepada berbagai pihak yang telah membantu, membimbing, dan memberikan dukungan selama proses penyusunan skripsi ini, antara lain:

\begin{enumerate}
    \item Nenek tercinta yang selalu mengunjungi dan membawakan makanan dengan penuh kasih sayang.
    \item Kakak dan adik yang selalu peduli dan memberikan semangat kepada penulis.
    \item Makci yang telah memberikan kasih sayang layaknya ibu kedua bagi penulis.
    \item Seluruh teman-teman seperjuangan yang telah memberikan dukungan dan kebersamaan selama masa perkuliahan.
    \item Dr. Puguh Wahyu Prasetyo, S.Si, M.Sc., selaku dosen pembimbing yang telah sabar membimbing, mengarahkan, dan memotivasi penulis hingga skripsi ini selesai.
    \item Ibu Rima Aksen Cahdriyana, M.Pd., yang telah membimbing penulis dalam berbagai kesempatan lomba dan memberikan ilmu yang bermanfaat.
    \item Seluruh dosen Program Studi Pendidikan Matematika Universitas Ahmad Dahlan yang telah memberikan ilmu dan pengalaman berharga selama penulis menempuh pendidikan.
    \item Bapak Wibowo Ramadhiyanto, S.Pd., selaku pendidik matematika di SMP Muhammadiyah 9 Yogyakarta, yang telah membantu penulis selama proses penelitian berlangsung.
    \item Jajaran Rektorat, Dekanat, dan Program Studi Pendidikan Matematika Universitas Ahmad Dahlan yang telah mengelola institusi dengan baik sehingga penulis dapat menyelesaikan studi.
\end{enumerate}

\hspace*{1cm}Penulis menyadari bahwa skripsi ini masih jauh dari kata sempurna. Oleh karena itu, kritik dan saran yang membangun sangat penulis harapkan demi perbaikan di masa mendatang. Semoga skripsi ini dapat memberikan manfaat bagi semua pihak yang membacanya, khususnya bagi pengembangan dunia pendidikan matematika.

\vspace{0.5cm}

\begin{flushright}
    \begin{minipage}{0.4\textwidth}
        \centering
        Yogyakarta, 16 Juni 2026 \\
        Penulis
    \end{minipage}
\end{flushright}

\vspace{0.5cm}

\begin{center}
    \textit{Wassalamu'alaikum warahmatullahi wabarakatuh.}
\end{center}